\section{Methodology}
\label{sec:method}

\para{Research question and hypotheses.} We ask whether having an LLM simulate the target reader
in its chain-of-thought roughly once per sentence produces writing that is (i) better fit to that
reader, (ii) higher in overall quality, and (iii) more genuinely adapted to the reader, compared
to plain generation, a token-matched generic reasoning control, and one-shot audience modeling. We
decompose this into five pre-registered hypotheses: \textbf{H1} (audience fit): \persentence
beats all baselines on a separate reader model's fit rating; \textbf{H2} (quality): \persentence
beats baselines on a cross-family quality judge; \textbf{H3} (adaptation): \persentence shows the
largest audience-adaptation gap; \textbf{H4} (not just tokens): \persentence beats generic
reasoning despite both adding thinking tokens; \textbf{H5} (diversity/cost): \persentence reduces
generic sameness at acceptable cost. The independent variable is the generation \emph{condition};
dependent variables are audience-fit score, pairwise quality win-rate, adaptation gap, diversity
metrics, and length/token cost.

\subsection{Task set}

We use \textbf{audience-targeted writing}, the sharpest test of audience modeling, comprising 60
tasks. \textbf{48 explanatory} tasks pair a concept with a specific reader: ``Explain
\{concept\} for \{audience\}'' in roughly 150 words. Twelve concepts span science, finance, and
technology (vaccines, compound interest, black holes, how an LLM generates text, antibiotic
resistance, blockchain, seasons, mRNA, the stock market, encryption, photosynthesis, greenhouse
warming), crossed with four audiences that vary in expertise and disposition: \textbf{an
8-year-old child}, \textbf{a non-technical executive}, \textbf{a PhD-level expert}, and \textbf{a
skeptic who distrusts the topic}. \textbf{12 creative} tasks draw prompts from
\writingprompts~\citep{fan2018writingprompts} with an assigned target readership, for breadth.

\subsection{Conditions}

The only manipulation is \emph{how the model reasons about the reader}; crucially, every
condition is given identical information about who the reader is, so any difference is
attributable to the reasoning strategy rather than to privileged audience information. Generation
uses a clean two-step protocol---a reasoning turn, then a ``write ONLY the final piece'' turn---so
the judged text never contains chain-of-thought scaffolding.

\begin{itemize}[leftmargin=*,itemsep=1pt,topsep=2pt]
    \item \plain: no reasoning turn---write directly (the ``annoying'' egocentric default).
    \item \cot: ``think step by step to make it excellent'' (generic reasoning;
    \textbf{token-matched control} that isolates audience simulation from added
    thinking tokens).
    \item \simtom: describe the reader once (background, knowledge, what they care about, tone),
    then write.
    \item \persentence \textit{(intervention)}: plan sentence by sentence; after each sentence,
    simulate how \emph{this reader} reacts (confused, bored, skeptical, curious), then choose the
    next sentence.
\end{itemize}

\subsection{Models}

We deliberately keep the generator and both judges in \textbf{different model families} to avoid
self-preference and shortcut inflation~\citep{shapira2023clever}. The \textbf{generator} is
\gptmini (primary) with \llama for replication. The \textbf{quality judge} is \claude, used for
blind, order-balanced pairwise comparison. The \textbf{audience/listener model} is \gemini, which
\emph{role-plays the target reader} to give absolute fit, understanding, and engagement ratings
($1$--$10$) and pairwise ``which serves this reader better'' judgments, following communication-based
evaluation~\citep{fried2023pragmatics}. All calls run through OpenRouter with a fixed seed of 42.

\subsection{Evaluation metrics}

\begin{enumerate}[leftmargin=*,itemsep=1pt,topsep=2pt]
    \item \textbf{Audience fit / understanding / engagement}: the audience model, role-playing the
    reader, rates each piece $1$--$10$.
    \item \textbf{Overall writing quality}: the cross-family judge compares pairs blind, in both
    orders, controlling for position bias; we report win-rate (0.5 = tie).
    \item \textbf{Adaptation gap} (the H3 control): each concept is written for a \emph{child} and
    for an \emph{expert}; both readers rate both pieces. The gap is $\text{matched} -
    \text{crossed}$ fit. A method that truly tailors to the reader scores high for its intended
    reader and low for the other, yielding a large gap; a generic-quality method yields a gap near
    zero. This isolates real audience use from a generic quality boost.
    \item \textbf{Homogenization}: distinct-2 and self-BLEU computed within each condition.
    \item \textbf{Cost}: output length in words and reasoning (``thinking'') characters.
\end{enumerate}

\subsection{Protocol and reproducibility}

Generation uses temperature $0.7$; judging uses temperature $0.0$; the seed is $42$ throughout.
Every model call is disk-cached by a hash of (model, messages, params), so the pipeline is fully
reproducible and re-runs are free. In total we ran $240$ primary generations, $192$ replication
generations, and roughly $1{,}400$ judge calls, yielding $3{,}876$ cached responses at an estimated
API cost of \$8--15. For statistics we use paired designs across tasks: Wilcoxon signed-rank tests
on paired differences, bootstrap $95\%$ confidence intervals for win-rates (against a null of
$0.5$) and mean differences, Cohen's $d$ / rank-biserial effect sizes, and
Holm--Bonferroni correction~\citep{holm1979simple} across the multiple condition comparisons, at
$\alpha = 0.05$. An initial marker-based text extractor leaked audience-analysis preamble into
about $23\%$ of \simtom pieces; we replaced it with the two-step protocol (zero leaks) and re-ran
everything, so all reported numbers come from the clean pipeline.
